%%
%% part2-book-components/ch11-graphics.tex
%%
%% Chapter 11: Graphics and Diagrams — Function plots,
%% statistical charts, flowcharts, trees, graphs, 3D
%% surfaces, pie charts, and timelines.
%%

\chapter{نمودار و شکل}
\label{chap:graphics}

\section{چرا این موضوع مهم است؟}
\label{sec:graphics-why}

در کتاب‌های فنی و علمی، نمودارها و شکل‌ها نقش
محوری در انتقال مفاهیم پیچیده دارند. یک نمودار
حرفه‌ای می‌تواند جای چند صفحه متن را بگیرد و
درک مطلب را چند برابر کند.

\lr{MatinBook} از \lr{TikZ}، \lr{PGFPlots} و
\lr{pgf-pie} برای ترسیم نمودارها و شکل‌های
حرفه‌ای استفاده می‌کند. در این فصل، با انواع
نمودارها و روش ساخت آن‌ها آشنا می‌شوید.

\mnote{\lr{TikZ} یک زبان برنامه‌نویسی برای ترسیم شکل‌های برداری است. \lr{PGFPlots} نیز بر پایه‌ی \lr{TikZ} برای رسم نمودارهای داده‌ای ساخته شده است.}

\section{ساختار پایه}
\label{sec:graphics-basic}

هر شکل در \lr{MatinBook} از یک ساختار پایه
پیروی می‌کند:

\begin{itemize}
    \item محیط \env{figure} برای نگهداری شکل.
    \item \cmd{centering} برای وسط‌چین کردن.
    \item محیط \env{latin} برای حفظ جهت لاتین.
    \item محیط \env{tikzpicture} یا
    \env{axis} برای ترسیم.
    \item \cmd{caption} برای کپشن فارسی.
    \item \cmd{label} برای ارجاع.
\end{itemize}

\mnote{محیط \lr{tikzpicture} باید همیشه داخل \lr{\textbackslash begin\{latin\}} قرار گیرد تا جهت ترسیم درست حفظ شود.}

\subsection{نمونه‌ی پایه}
\label{sec:graphics-basic-example}

\begin{latin}
\begin{minted}{latex}
\begin{figure}[h]
\centering
\begin{latin}
\begin{tikzpicture}
    \draw[matinblue, thick] (0,0) -- (3,2);
    \draw[matinred, thick] (0,2) -- (3,0);
\end{tikzpicture}
\end{latin}
\caption{|\pc{یک شکل ساده}|}
\label{fig:simple}
\end{figure}
\end{minted}
\end{latin}

\mnote{دستور \lr{[h]} به \lr{LaTeX} می‌گوید که شکل را «همین‌جا» قرار دهد. اگر فضا کافی نباشد، به صفحه‌ی بعد می‌رود.}

\section{نمودار توابع}
\label{sec:graphics-functions}

پرکاربردترین نوع نمودار در کتاب‌های فنی،
نمودار توابع ریاضی است. \lr{MatinBook} از
\lr{PGFPlots} برای این کار استفاده می‌کند.

\subsection{نمودار تابع ساده}
\label{sec:graphics-simple-function}

\begin{latin}
\begin{minted}{latex}
\begin{figure}[h]
\centering
\begin{latin}
\begin{tikzpicture}
  \begin{axis}[
    width=\textwidth, height=6cm,
    xlabel={$x$},
    ylabel={$f(x)$},
    grid=major,
    domain=-3:3,
    samples=100
  ]
    \addplot[matinblue, thick] {x^2};
    \addlegendentry{$x^2$}
  \end{axis}
\end{tikzpicture}
\end{latin}
\caption{|\pc{نمودار تابع}|\lr{$f(x) = x^{2}$}}
\label{fig:parabola}
\end{figure}
\end{minted}
\end{latin}

خروجی:

\begin{figure}[h]
\centering
\begin{latin}
\begin{tikzpicture}
  \begin{axis}[
    width=0.8\textwidth, height=5cm,
    xlabel={$x$},
    ylabel={$f(x)$},
    grid=major,
    domain=-3:3,
    samples=100
  ]
    \addplot[matinblue, thick] {x^2};
    \addlegendentry{$x^2$}
  \end{axis}
\end{tikzpicture}
\end{latin}
\caption{نمودار تابع \lr{$f(x) = x^{2}$}}
\label{fig:parabola-ch11}
\end{figure}

\mnote{مقدار \lr{width=\textbackslash textwidth} باعث می‌شود نمودار دقیقاً به عرض متن چسبیده و از صفحه سرریز نکند.}

\subsection{نمودار چند تابع}
\label{sec:graphics-multiple-functions}

\begin{latin}
\begin{minted}{latex}
\begin{axis}[
    width=\textwidth, height=6cm,
    xlabel={$x$},
    ylabel={$f(x)$},
    grid=major,
    legend pos=north west,
    domain=-2:2,
    samples=100
]
    \addplot[matinblue, thick] {x^2};
    \addlegendentry{$x^2$}

    \addplot[matinred, thick] {x^3};
    \addlegendentry{$x^3$}

    \addplot[matinorange, thick] {exp(x)};
    \addlegendentry{$e^x$}
\end{axis}
\end{minted}
\end{latin}

خروجی:

\begin{figure}[h]
\centering
\begin{latin}
\begin{tikzpicture}
  \begin{axis}[
    width=0.8\textwidth, height=5cm,
    xlabel={$x$},
    ylabel={$f(x)$},
    grid=major,
    legend pos=north west,
    domain=-2:2,
    samples=100
  ]
    \addplot[matinblue, thick] {x^2};
    \addlegendentry{$x^2$}
    \addplot[matinred, thick] {x^3};
    \addlegendentry{$x^3$}
    \addplot[matinorange, thick] {exp(x)};
    \addlegendentry{$e^x$}
  \end{axis}
\end{tikzpicture}
\end{latin}
\caption{مقایسه‌ی سه تابع}
\label{fig:three-functions}
\end{figure}

\mnote{پالت رنگ \lr{matinbook} به‌طور خودکار برای نمودارهای چندسری اعمال می‌شود. می‌توانید رنگ‌ها را دستی هم مشخص کنید.}

\section{نمودارهای آماری}
\label{sec:graphics-statistical}

\subsection{نمودار میله‌ای}
\label{sec:graphics-bar}

\begin{latin}
\begin{minted}{latex}
\begin{axis}[
    width=\textwidth, height=6cm,
    ybar,
    symbolic x coords={A, B, C, D, E},
    xtick=data,
    ylabel={|\pc{مقدار}|},
    bar width=15pt,
    nodes near coords
]
    \addplot[fill=matinblue!70] coordinates {
        (A, 25) (B, 40) (C, 30) (D, 50) (E, 35)
    };
\end{axis}
\end{minted}
\end{latin}

\mnote{در \lr{PGFPlots}، برچسب‌های محور باید لاتین باشند. برای معادل فارسی، از کپشن یا متن استفاده کنید.}

\subsection{نمودار خطی}
\label{sec:graphics-line}

\begin{latin}
\begin{minted}{latex}
\begin{axis}[
    width=\textwidth, height=6cm,
    xlabel={$n$},
    ylabel={Operations},
    grid=major,
    legend pos=north west,
    domain=0:50,
    samples=100
]
    \addplot[matinred, thick] {x^2};
    \addlegendentry{$O(n^2)$}

    \addplot[matinblue, thick] {x*ln(x)};
    \addlegendentry{$O(n \log n)$}

    \addplot[matingreen, thick] {x};
    \addlegendentry{$O(n)$}
\end{axis}
\end{minted}
\end{latin}

\subsection{نمودار پراکندگی}
\label{sec:graphics-scatter}

\begin{latin}
\begin{minted}{latex}
\begin{axis}[
    width=\textwidth, height=6cm,
    xlabel={X},
    ylabel={Y},
    grid=major
]
    \addplot[only marks, mark=*, color=matinblue] coordinates {
        (1, 45) (2, 52) (3, 58) (4, 65) (5, 70)
    };
\end{axis}
\end{minted}
\end{latin}

\mnote{برای نمودار پراکندگی، از \lr{only marks} استفاده کنید تا نقاط به هم متصل نشوند.}

\section{فلوچارت‌ها}
\label{sec:graphics-flowcharts}

فلوچارت‌ها برای نمایش فرآیندها و الگوریتم‌ها
استفاده می‌شوند. \lr{MatinBook} چند استایل
آماده برای فلوچارت ارائه می‌دهد:

\begin{itemize}
    \item \lr{block}: مستطیل با گوشه‌های گرد.
    \item \lr{decision}: لوزی برای تصمیم.
    \item \lr{arrow}: فلش ضخیم.
\end{itemize}

\begin{latin}
\begin{minted}{latex}
\begin{tikzpicture}[node distance=2cm]
    \node[block] (start) {Start};
    \node[decision, below of=start] (check) {$x > 0$?};
    \node[block, below of=check] (positive) {Positive};
    \node[block, right of=check, node distance=3cm] (negative) {Negative};

    \draw[arrow] (start) -- (check);
    \draw[arrow] (check) -- node[left] {Yes} (positive);
    \draw[arrow] (check) -- node[above] {No} (negative);
\end{tikzpicture}
\end{minted}
\end{latin}

خروجی:

\begin{figure}[h]
\centering
\begin{latin}
\begin{tikzpicture}[node distance=1.5cm, scale=0.9, transform shape]
    \node[block] (start) {Start};
    \node[decision, below of=start] (check) {$x > 0$?};
    \node[block, below of=check] (positive) {Positive};
    \node[block, right of=check, node distance=3cm] (negative) {Negative};

    \draw[arrow] (start) -- (check);
    \draw[arrow] (check) -- node[left] {Yes} (positive);
    \draw[arrow] (check) -- node[above] {No} (negative);
\end{tikzpicture}
\end{latin}
\caption{فلوچارت یک الگوریتم ساده}
\label{fig:flowchart-ch11}
\end{figure}

\mnote{در فلوچارت‌ها، متن داخل گره‌ها باید لاتین باشد. برای متن فارسی، از \lr{\textbackslash rl\{...\}} استفاده کنید.}

\section{درخت‌ها}
\label{sec:graphics-trees}

درخت‌ها برای نمایش ساختارهای سلسله‌مراتبی
استفاده می‌شوند. \lr{MatinBook} یک استایل
\lr{treenode} برای گره‌های درخت ارائه می‌دهد:

\begin{latin}
\begin{minted}{latex}
\begin{tikzpicture}[
    level distance=1.5cm,
    level 1/.style={sibling distance=3cm},
    level 2/.style={sibling distance=1.5cm}
]
    \node[treenode] {8}
        child {node[treenode] {3}
            child {node[treenode] {1}}
            child {node[treenode] {6}}
        }
        child {node[treenode] {10}
            child {node[treenode] {9}}
            child {node[treenode] {14}}
        };
\end{tikzpicture}
\end{minted}
\end{latin}

خروجی:

\begin{figure}[h]
\centering
\begin{latin}
\begin{tikzpicture}[
    level distance=1.5cm,
    scale=0.9,
    transform shape,
    level 1/.style={sibling distance=3cm},
    level 2/.style={sibling distance=1.5cm}
]
    \node[treenode] {8}
        child {node[treenode] {3}
            child {node[treenode] {1}}
            child {node[treenode] {6}}
        }
        child {node[treenode] {10}
            child {node[treenode] {9}}
            child {node[treenode] {14}}
        };
\end{tikzpicture}
\end{latin}
\caption{درخت جستجوی دودویی}
\label{fig:bst-ch11}
\end{figure}

\section{گراف‌ها}
\label{sec:graphics-graphs}

گراف‌ها برای نمایش روابط بین موجودیت‌ها استفاده
می‌شوند. \lr{MatinBook} یک استایل \lr{vertex}
برای رأس‌های گراف ارائه می‌دهد:

\begin{latin}
\begin{minted}{latex}
\begin{tikzpicture}[
    vertex/.style={circle, draw=matindarkblue,
                   fill=matinlightblue!30, minimum size=1cm}
]
    \node[vertex] (A) at (0,0) {A};
    \node[vertex] (B) at (3,0) {B};
    \node[vertex] (C) at (1.5,-2) {C};

    \draw[thick] (A) -- (B);
    \draw[thick] (B) -- (C);
    \draw[thick] (A) -- (C);
\end{tikzpicture}
\end{minted}
\end{latin}

\mnote{برای گراف‌های پیچیده، از کتابخانه‌ی \lr{graphs} در \lr{TikZ} استفاده کنید که امکان تعریف خودکار گراف‌ها را فراهم می‌کند.}

\section{نمودار سه‌بعدی}
\label{sec:graphics-3d}

\lr{PGFPlots} از نمودارهای سه‌بعدی نیز پشتیبانی
می‌کند. برای این کار، باید کتابخانه‌ی
\lr{colormaps} را فعال کنید:

\begin{latin}
\begin{minted}{latex}
\usepgfplotslibrary{colormaps}

\begin{axis}[
    width=\textwidth, height=8cm,
    view={45}{35},
    xlabel={$x$}, ylabel={$y$}, zlabel={$f$}
]
    \addplot3[
        surf, domain=-5:5, domain y=-5:5,
        samples=40, samples y=40,
        colormap/viridis
    ] {sin(deg(sqrt(x^2+y^2)))};
\end{axis}
\end{minted}
\end{latin}

\mnote{نمودارهای سه‌بعدی ممکن است در بعضی از نسخه‌های \lr{PGFPlots} کند باشند. صبور باشید.}

\section{نمودار دایره‌ای}
\label{sec:graphics-pie}

\lr{MatinBook} از \lr{pgf-pie} برای نمودار
دایره‌ای پشتیبانی می‌کند:

\begin{latin}
\begin{minted}{latex}
\begin{tikzpicture}
    \pie[
        radius=3,
        text=legend,
        color={matinblue, matingreen, matinorange, matinred}
    ]{
        35/Python,
        25/JavaScript,
        20/Java,
        12/C++,
        8/Other
    }
\end{tikzpicture}
\end{minted}
\end{latin}

\mnote{در نمودار دایره‌ای، برچسب‌ها باید لاتین باشند. برای متن فارسی، از کپشن استفاده کنید.}

\section{خط زمانی}
\label{sec:graphics-timeline}

خط زمانی برای نمایش رویدادها در طول زمان
استفاده می‌شود:

\begin{latin}
\begin{minted}{latex}
\begin{tikzpicture}[xscale=0.7]
    \draw[thick,->] (0,0) -- (14,0);

    \foreach \x/\year in {0/1990, 3.5/2000, 7/2010, 10.5/2020, 14/2025} {
        \draw[thick] (\x,-0.2) -- (\x,0.2);
        \node[below] at (\x,-0.3) {\small\year};
    }

    \node[event, above] at (0,1.2) {Python\\Created};
    \node[event, above] at (7,1.2) {Rust\\First Stable};
\end{tikzpicture}
\end{minted}
\end{latin}

\section{فاصله‌گذاری و اندازه}
\label{sec:graphics-sizing}

یکی از مشکلات رایج در نمودارها، سرریز به
حاشیه است. برای جلوگیری از این مشکل:

\begin{itemize}
    \item از \lr{width=\textbackslash textwidth} برای
    \env{axis} استفاده کنید.
    \item از \lr{scale=0.9, transform shape}
    برای \env{tikzpicture} استفاده کنید.
    \item عرض ثابت (مثل \lr{width=14cm}) را
    حذف کنید.
\end{itemize}

\mnote{عرض \lr{\textbackslash textwidth} در متن اصلی حدود \lr{۱۰.۳} سانتی‌متر است. اگر عرض نمودار بیشتر باشد، سرریز می‌کند.}

\section{ارجاع به شکل}
\label{sec:graphics-references}

برای ارجاع به شکل، از \cmd{cref} استفاده کنید:

\begin{latin}
\begin{minted}{latex}
|\pc{نمودار}|\cref{fig:parabola}|\pc{تابع درجه دو را نشان می‌دهد.}|%
\end{minted}
\end{latin}

خروجی: نمودار \cref{fig:parabola-ch11} تابع
درجه دو را نشان می‌دهد.

\mnote{نام‌گذاری برچسب شکل‌ها با پیشوند \lr{fig:} انجام می‌شود: \lr{fig:parabola}.}

\section{متن فارسی در \lr{tikzpicture}}
\label{sec:graphics-persian}

برای متن فارسی داخل \env{tikzpicture}، باید از
\cmd{rl} استفاده کنید:

\begin{latin}
\begin{minted}{latex}
\begin{tikzpicture}
    \node[block] (start) {|\rl{شروع}|};
    \node[block, below of=start] (end) {|\rl{پایان}|};
    \draw[arrow] (start) -- (end);
\end{tikzpicture}
\end{minted}
\end{latin}

\mnote{فراموش کردن \lr{\textbackslash rl} باعث می‌شود متن فارسی به‌صورت چپ‌چین و نامفهوم نمایش داده شود.}

\section{خطاهای رایج}
\label{sec:graphics-mistakes}

\subsection{فراموش کردن \env{latin}}
\label{sec:graphics-mistake-latin}

اگر \env{tikzpicture} را داخل \env{latin} قرار
ندهید، ممکن است جهت ترسیم به‌هم بریزد و
شکل در جای اشتباه ظاهر شود.

\textbf{راه‌حل:} همیشه \env{tikzpicture} را داخل
\env{latin} قرار دهید.

\subsection{سرریز نمودار به حاشیه}
\label{sec:graphics-mistake-overflow}

اگر عرض نمودار بیشتر از \cmd{textwidth} باشد،
به حاشیه سرریز می‌کند.

\textbf{راه‌حل:} از \lr{width=\textbackslash textwidth}
استفاده کنید، یا مقیاس را کاهش دهید.

\subsection{فراموش کردن \cmd{rl} در \lr{tikzpicture}}
\label{sec:graphics-mistake-rl}

اگر متن فارسی را داخل \env{tikzpicture} بدون
\cmd{rl} بنویسید، به‌صورت چپ‌چین و نامفهوم
نمایش داده می‌شود.

\textbf{راه‌حل:} متن فارسی را در \lr{\textbackslash rl\{...\}}
قرار دهید.

\section{تمرین‌ها}
\label{sec:graphics-exercises}

\begin{exercise}
    یک نمودار تابع ساده بسازید که عرض آن
    دقیقاً به عرض متن چسبیده باشد.
    \label{exc:graphics-simple}
\end{exercise}

\begin{exercise}
    یک فلوچارت با سه گره بسازید و متن‌های
    فارسی آن را با \cmd{rl} پیچیده کنید.
    \label{exc:graphics-flowchart}
\end{exercise}

\begin{exercise}
    یک درخت دودویی با سه سطح بسازید.
    \label{exc:graphics-tree}
\end{exercise}

\begin{exercise}
    یک نمودار دایره‌ای با چهار بخش بسازید.
    \label{exc:graphics-pie}
\end{exercise}

\section{خلاصه}
\label{sec:graphics-summary}

در این فصل، با نمودار و شکل در \lr{MatinBook}
آشنا شدیم:

\begin{itemize}
    \item \lr{MatinBook} از \lr{TikZ}، \lr{PGFPlots}
    و \lr{pgf-pie} برای ترسیم استفاده می‌کند.

    \item نمودار توابع با \env{axis} و
    \cmd{addplot} ساخته می‌شوند.

    \item نمودارهای آماری شامل میله‌ای، خطی و
    پراکندگی هستند.

    \item فلوچارت‌ها با استایل‌های \lr{block}،
    \lr{decision} و \lr{arrow} ساخته می‌شوند.

    \item درخت‌ها با استایل \lr{treenode} ساخته
    می‌شوند.

    \item گراف‌ها با استایل \lr{vertex} ساخته
    می‌شوند.

    \item \env{tikzpicture} باید همیشه داخل
    \env{latin} قرار گیرد.

    \item متن فارسی داخل \env{tikzpicture} باید
    با \cmd{rl} پیچیده شود.

    \item خطاهای رایج شامل فراموش کردن
    \env{latin}، سرریز نمودار، و فراموش کردن
    \cmd{rl} است.
\end{itemize}

در فصل بعدی (\cref{chap:tables})، با جدول
در \lr{MatinBook} آشنا می‌شوید و یاد می‌گیرید
که چگونه جدول‌های حرفه‌ای بسازید.